% =============================================================================
%  PAES Studio: presentación de repaso (plantilla Beamer 16:9)
%  Lista para Overleaf: Nuevo proyecto → Subir proyecto → compilar con pdfLaTeX.
%
%  Estética minimalista con la paleta Zinc de la app: sin degradados, sin
%  símbolos de navegación, líneas finas y alternativas con su tecla A–E.
%  Comandos:
%    \alternativa{A}{texto}       alternativa
%    \alternativaok{C}{texto}     alternativa correcta (para la diapositiva de solución)
%    \paso{n}{texto}      paso numerado del desarrollo
%    \ejebarra{Eje}{0.42}{42 %}   barra proporcional de frecuencia
%  Entornos: cajaclave, cajatrampa, lineatiempo (filas «fecha & hecho \\»).
%
%  El generador scripts/pipeline/generate_latex_summary.py reemplaza los bloques
%  delimitados por %%<NOMBRE>%% … %%</NOMBRE>%%.
% =============================================================================
\documentclass[aspectratio=169,11pt]{beamer}

\usepackage{lmodern}        % antes de fontenc/babel: no depende de las fuentes EC
\usepackage[T1]{fontenc}
\usepackage[utf8]{inputenc}
\usepackage[spanish,es-noshorthands,es-tabla]{babel}
\spanishplainpercent        % \% sin espaciado automático (evita choques en modo matemático)
\usepackage{textcomp}
\usepackage{amsmath,amssymb}
\usepackage{booktabs}
\usepackage{array}
\usepackage[most]{tcolorbox}

% ── Paleta Zinc ──────────────────────────────────────────────────────────────
\definecolor{zinc950}{HTML}{09090B}
\definecolor{zinc900}{HTML}{18181B}
\definecolor{zinc700}{HTML}{3F3F46}
\definecolor{zinc500}{HTML}{71717A}
\definecolor{zinc300}{HTML}{D4D4D8}
\definecolor{zinc200}{HTML}{E4E4E7}
\definecolor{zinc50}{HTML}{FAFAFA}
\definecolor{esmeralda}{HTML}{047857}
\definecolor{rosa}{HTML}{BE123C}

\newcommand{\sen}{\operatorname{sen}}

% ── Tema mínimo ──────────────────────────────────────────────────────────────
\usefonttheme{professionalfonts}
\setbeamertemplate{navigation symbols}{}
\setbeamercolor{background canvas}{bg=white}
\setbeamercolor{normal text}{fg=zinc900}
\setbeamercolor{structure}{fg=zinc900}
\setbeamercolor{alerted text}{fg=zinc950}
\setbeamercolor{item}{fg=zinc500}
\setbeamerfont{frametitle}{series=\bfseries,size=\large}
\setbeamertemplate{itemize items}{\raisebox{0.15ex}{\color{zinc500}\scriptsize$\bullet$}}
\setbeamertemplate{enumerate items}{\color{zinc500}\insertenumlabel.}
\setbeamersize{text margin left=1.6cm,text margin right=1.6cm}

\makeatletter
\setbeamertemplate{frametitle}{%
  \vspace*{0.9em}%
  {\usebeamerfont{frametitle}\color{zinc950}\insertframetitle}\par
  \ifx\insertframesubtitle\@empty\else{\small\color{zinc500}\insertframesubtitle}\par\fi
  \vspace{0.25em}{\color{zinc200}\rule{\textwidth}{0.4pt}}%
}
\makeatother
\setbeamertemplate{footline}{%
  \hspace*{1.6cm}{\color{zinc500}\tiny\insertshorttitle}\hfill
  {\color{zinc500}\tiny\ttfamily\insertframenumber/\inserttotalframenumber}\hspace*{1.6cm}\par\vspace*{0.7em}%
}
\setbeamertemplate{title page}{%
  \vfill
  {\color{zinc500}\small PAES Studio}\par\vspace{0.6em}
  {\usebeamerfont{title}\bfseries\Huge\color{zinc950}\inserttitle}\par\vspace{0.5em}
  {\large\color{zinc700}\insertsubtitle}\par\vspace{1.2em}
  {\color{zinc300}\rule{0.35\textwidth}{0.4pt}}\par\vspace{0.6em}
  {\small\color{zinc500}\insertdate}\par
  \vfill
}

% ── Componentes ──────────────────────────────────────────────────────────────
\newtcbox{\tecla}{on line, boxrule=0.4pt, arc=1.5pt, colback=zinc50, colframe=zinc300,
  left=3pt, right=3pt, top=1pt, bottom=1pt, fontupper=\ttfamily\footnotesize\color{zinc700}}
\newtcbox{\teclaok}{on line, boxrule=0.4pt, arc=1.5pt, colback=esmeralda, colframe=esmeralda,
  left=3pt, right=3pt, top=1pt, bottom=1pt, fontupper=\ttfamily\footnotesize\bfseries\color{white}}
\newtcolorbox{opcion}[1][]{enhanced, boxrule=0.4pt, arc=2pt, colback=white, colframe=zinc300,
  left=6pt, right=6pt, top=3pt, bottom=3pt, before skip=3pt, after skip=3pt, #1}
\newcommand{\alternativa}[2]{\begin{opcion}\tecla{#1}\hspace{0.6em}#2\end{opcion}}
\newcommand{\alternativaok}[2]{\begin{opcion}[colframe=esmeralda!60, colback=esmeralda!5]\teclaok{#1}\hspace{0.6em}#2\end{opcion}}
\newcommand{\paso}[2]{\par\noindent\makebox[1.4em][l]{\color{zinc500}\ttfamily\small #1}\parbox[t]{\dimexpr\linewidth-1.4em}{#2}\par\vspace{0.45em}}
\newtcolorbox{cajaclave}[1][]{enhanced, boxrule=0.4pt, arc=2pt, colback=zinc50, colframe=zinc300,
  borderline west={2pt}{0pt}{zinc900}, left=7pt, right=7pt, top=3pt, bottom=4pt,
  fonttitle=\bfseries\small, coltitle=zinc900, colbacktitle=zinc50, titlerule=0pt, title={#1}}
\newtcolorbox{cajatrampa}[1][Trampas frecuentes]{enhanced, boxrule=0.4pt, arc=2pt, colback=white, colframe=rosa!35,
  borderline west={2pt}{0pt}{rosa}, left=7pt, right=7pt, top=3pt, bottom=4pt,
  fonttitle=\bfseries\small, coltitle=zinc900, colbacktitle=white, titlerule=0pt, title={#1}}
\newenvironment{lineatiempo}
  {\footnotesize\renewcommand{\arraystretch}{1.15}\begin{tabular}{@{}>{\leavevmode\bfseries\color{zinc700}}p{0.14\textwidth}p{0.74\textwidth}@{}}}
  {\end{tabular}}
\newlength{\anchobarra}
\newcommand{\ejebarra}[3]{%
  \par\noindent\makebox[0.34\linewidth][l]{\small #1}%
  \setlength{\anchobarra}{0.5\linewidth}%
  \makebox[\anchobarra][l]{{\color{zinc200}\rule[0.2ex]{\anchobarra}{1.1ex}}\hspace*{-\anchobarra}{\color{zinc900}\rule[0.2ex]{#2\anchobarra}{1.1ex}}}%
  \hspace{0.8em}{\small\ttfamily #3}\par\vspace{0.5em}}

%%<META>%%
\title[Repaso PAES: Competencia Matemática 2]{Repaso PAES}
\subtitle{Competencia Matemática 2 (M2)}
\date{8 de noviembre de 2026}
%%</META>%%

\begin{document}

\begin{frame}[plain]
  \titlepage
\end{frame}

%%<CONTENIDO>%%
\begin{frame}{Contenidos}

  \framesubtitle{5 ejes, en el orden del temario}

  \begin{enumerate}\small

    \item \textbf{Números}: Números reales y racionalización, Logaritmos, Matemática financiera

    \item \textbf{Álgebra y funciones}: Sistemas de ecuaciones: tipos de solución, Función potencia, exponencial y logarítmica, Funciones trigonométricas seno y coseno, Ecuación y función cuadrática (contenido de M1)

    \item \textbf{Geometría}: Homotecia de figuras planas, Razones trigonométricas, Relaciones métricas en la circunferencia, Esfera: superficie y volumen, Rectas en el plano

    \item \textbf{Probabilidad y estadística}: Medidas de dispersión, Probabilidad condicional, Permutación y combinatoria, Modelos binomial y normal

    \item \textbf{Habilidades transversales}: Evaluación de suficiencia de datos

  \end{enumerate}

\end{frame}

\begin{frame}{Qué evalúa la prueba}

  \framesubtitle{Distribución por eje en 8 ensayos (363 preguntas)}

  \vspace{0.8em}

  \ejebarra{Álgebra y funciones}{0.284}{28 \%}

  \ejebarra{Números}{0.264}{26 \%}

  \ejebarra{Geometría}{0.251}{25 \%}

  \ejebarra{Probabilidad y estadística}{0.201}{20 \%}

\end{frame}

\section{Números}

\begin{frame}[plain]
  \vfill
  {\color{zinc500}\small Eje}\par\vspace{0.3em}{\Large\bfseries Números}\par\vspace{0.6em}
  {\color{zinc700} Números reales, logaritmos y matemática financiera.}\par\vspace{1em}
  {\small\color{zinc500} Números reales y racionalización, Logaritmos, Matemática financiera}
  \vfill
\end{frame}

\begin{frame}[allowframebreaks]{Números reales y racionalización}
  \framesubtitle{Teoría esencial}
  \begin{itemize}\small
    \item $\sqrt{2}\cdot\sqrt{8}=\sqrt{16}=4$ es racional: el producto de irracionales puede ser racional.
    \item El conjugado de $\sqrt{5}-1$ es $\sqrt{5}+1$.
    \item Para ordenar raíces, compara sus cuadrados.
  \end{itemize}
  \begin{cajaclave}[Racionalización]\footnotesize
\begin{gather*}
\frac{k}{\sqrt{a}}=\frac{k\sqrt{a}}{a} \\
\frac{k}{\sqrt{a}-\sqrt{b}}=\frac{k(\sqrt{a}+\sqrt{b})}{a-b}
\end{gather*}
  \end{cajaclave}
\end{frame}

\begin{frame}{Desafío}
  \framesubtitle{Números reales y racionalización}
  {\small ¿Cuál de las siguientes expresiones es igual a $\dfrac{6}{\sqrt{3}}$?}
  \vspace{0.4em}
  \alternativa{A}{$2\sqrt{3}$}
  \alternativa{B}{$3\sqrt{2}$}
  \alternativa{C}{$\sqrt{2}$}
  \alternativa{D}{$6\sqrt{3}$}
\end{frame}

\begin{frame}[allowframebreaks]{Solución}
  \framesubtitle{Amplifica por $\sqrt{3}$ para eliminar la raíz del denominador.}
  {\small
  \paso{1}{$\dfrac{6}{\sqrt{3}}\cdot\dfrac{\sqrt{3}}{\sqrt{3}}=\dfrac{6\sqrt{3}}{3}$.}
  \paso{2}{Simplifica: $2\sqrt{3}$.}
  }
  \alternativaok{A}{$2\sqrt{3}$}
  \begin{cajatrampa}\footnotesize
    \begin{itemize}
    \item Amplificar solo el denominador.
    \item Suponer que la suma o el producto de irracionales siempre es irracional.
    \end{itemize}
  \end{cajatrampa}
\end{frame}

\begin{frame}[allowframebreaks]{Logaritmos}
  \framesubtitle{Teoría esencial}
  \begin{itemize}\small
    \item $\log_2 32=5$ porque $2^5=32$.
    \item Si el argumento está entre 0 y 1, el logaritmo es negativo.
    \item $\log(x+y)\neq\log x+\log y$.
  \end{itemize}
  \begin{cajaclave}[Propiedades]\footnotesize
\begin{gather*}
\log_b(xy)=\log_b x+\log_b y \\
\log_b\frac{x}{y}=\log_b x-\log_b y \\
\log_b x^n=n\log_b x
\end{gather*}

\begin{gather*}
\log_b b=1 \\
\log_b 1=0 \\
\log_b a=\frac{\log a}{\log b}
\end{gather*}
  \end{cajaclave}
\end{frame}

\begin{frame}{Desafío}
  \framesubtitle{Logaritmos}
  {\small Si $\log 2\approx 0{,}301$, ¿cuál es el valor aproximado de $\log 50$?}
  \vspace{0.4em}
  \alternativa{A}{$0{,}602$}
  \alternativa{B}{$1{,}398$}
  \alternativa{C}{$1{,}699$}
  \alternativa{D}{$2{,}301$}
\end{frame}

\begin{frame}[allowframebreaks]{Solución}
  \framesubtitle{$50=\dfrac{100}{2}$, y $\log 100=2$.}
  {\small
  \paso{1}{$\log 50=\log 100-\log 2$.}
  \paso{2}{$2-0{,}301=1{,}699$.}
  }
  \alternativaok{C}{$1{,}699$}
  \begin{cajatrampa}\footnotesize
    \begin{itemize}
    \item Distribuir el logaritmo sobre una suma.
    \item Olvidar que el argumento debe ser positivo.
    \end{itemize}
  \end{cajatrampa}
\end{frame}

\begin{frame}[allowframebreaks]{Matemática financiera}
  \framesubtitle{Teoría esencial}
  \begin{itemize}\small
    \item Con 2 \% mensual, en 3 meses: simple $\cdot 1{,}06$; compuesto $\cdot 1{,}02^3\approx 1{,}0612$.
    \item El compuesto supera al simple desde el segundo período.
    \item Depreciación de $p\,\%$ anual: $V_n=V_0(1-p/100)^n$.
  \end{itemize}
  \begin{cajaclave}[Fórmulas]\footnotesize
\begin{gather*}
\text{simple: } C_n=C_0(1+n\,i) \\
\text{compuesto: } C_n=C_0(1+i)^n
\end{gather*}
  \end{cajaclave}
\end{frame}

\begin{frame}{Desafío}
  \framesubtitle{Matemática financiera}
  {\small Se invierten \$500 000 a un interés compuesto del 10 \% anual. ¿Cuánto interés se habrá ganado al cabo de 2 años?}
  \vspace{0.4em}
  \alternativa{A}{\$100 000}
  \alternativa{B}{\$105 000}
  \alternativa{C}{\$110 000}
  \alternativa{D}{\$605 000}
\end{frame}

\begin{frame}[allowframebreaks]{Solución}
  \framesubtitle{Calcula el monto final y réstale el capital.}
  {\small
  \paso{1}{$C_2=500\,000\cdot 1{,}1^2=500\,000\cdot 1{,}21=605\,000$.}
  \paso{2}{Interés ganado: $605\,000-500\,000=105\,000$.}
  }
  \alternativaok{B}{\$105 000}
  \begin{cajatrampa}\footnotesize
    \begin{itemize}
    \item Usar interés simple cuando el problema dice compuesto.
    \item Responder el monto final cuando se pide solo el interés ganado.
    \end{itemize}
  \end{cajatrampa}
\end{frame}

\section{Álgebra y funciones}

\begin{frame}[plain]
  \vfill
  {\color{zinc500}\small Eje}\par\vspace{0.3em}{\Large\bfseries Álgebra y funciones}\par\vspace{0.6em}
  {\color{zinc700} Sistemas de ecuaciones, funciones potencia, exponencial, logarítmica y trigonométricas, y la ecuación cuadrática de M1.}\par\vspace{1em}
  {\small\color{zinc500} Sistemas de ecuaciones: tipos de solución, Función potencia, exponencial y logarítmica, Funciones trigonométricas seno y coseno, Ecuación y función cuadrática (contenido de M1)}
  \vfill
\end{frame}

\begin{frame}[allowframebreaks]{Sistemas de ecuaciones: tipos de solución}
  \framesubtitle{Teoría esencial}
  \begin{itemize}\small
    \item Pendientes distintas: las rectas se cortan en un punto.
    \item Misma pendiente y distinto coeficiente de posición: rectas paralelas, sin solución.
    \item Una ecuación múltiplo de la otra: infinitas soluciones.
  \end{itemize}
  \begin{cajaclave}[Criterio]\footnotesize
\begin{gather*}
\frac{a}{d}\neq\frac{b}{e}\ \Rightarrow\ \text{única} \\
\frac{a}{d}=\frac{b}{e}\neq\frac{c}{f}\ \Rightarrow\ \text{ninguna} \\
\frac{a}{d}=\frac{b}{e}=\frac{c}{f}\ \Rightarrow\ \text{infinitas}
\end{gather*}
  \end{cajaclave}
\end{frame}

\begin{frame}{Desafío}
  \framesubtitle{Sistemas de ecuaciones: tipos de solución}
  {\small ¿Para qué valor de $k$ el sistema $2x+ky=5$ y $4x+6y=7$ no tiene solución?}
  \vspace{0.4em}
  \alternativa{A}{$k=2$}
  \alternativa{B}{$k=3$}
  \alternativa{C}{$k=6$}
  \alternativa{D}{$k=12$}
\end{frame}

\begin{frame}[allowframebreaks]{Solución}
  \framesubtitle{No hay solución cuando las rectas son paralelas y distintas.}
  {\small
  \paso{1}{$\dfrac{2}{4}=\dfrac{k}{6}\Rightarrow k=3$.}
  \paso{2}{Con $k=3$: $\dfrac{5}{7}\neq\dfrac{1}{2}$, así que son paralelas distintas.}
  }
  \alternativaok{B}{$k=3$}
  \begin{cajatrampa}\footnotesize
    \begin{itemize}
    \item Concluir «sin solución» solo porque las pendientes son iguales, sin revisar si las rectas coinciden.
    \item Comparar coeficientes sin ordenar antes las ecuaciones.
    \end{itemize}
  \end{cajatrampa}
\end{frame}

\begin{frame}[allowframebreaks]{Función potencia, exponencial y logarítmica}
  \framesubtitle{Teoría esencial}
  \begin{itemize}\small
    \item $f(0)=a$: el valor inicial.
    \item Crecimiento de $r\,\%$ por período: $b=1+\frac{r}{100}$.
    \item $\log_b x$ solo existe para $x>0$.
  \end{itemize}
  \begin{cajaclave}[Modelos]\footnotesize
\begin{gather*}
N(t)=N_0\cdot 2^{t/T}\ \text{(duplicación cada } T\text{)} \\
N(t)=N_0\left(\tfrac{1}{2}\right)^{t/T}\ \text{(vida media } T\text{)}
\end{gather*}
  \end{cajaclave}
\end{frame}

\begin{frame}{Desafío}
  \framesubtitle{Función potencia, exponencial y logarítmica}
  {\small Una sustancia tiene una vida media de 5 años. Si hoy hay 80 g, ¿cuánta quedará en 15 años?}
  \vspace{0.4em}
  \alternativa{A}{40 g}
  \alternativa{B}{20 g}
  \alternativa{C}{10 g}
  \alternativa{D}{5 g}
\end{frame}

\begin{frame}[allowframebreaks]{Solución}
  \framesubtitle{Cada 5 años la cantidad se reduce a la mitad.}
  {\small
  \paso{1}{En 15 años hay $\frac{15}{5}=3$ vidas medias.}
  \paso{2}{$80\cdot\left(\tfrac{1}{2}\right)^3=10$ g.}
  }
  \alternativaok{C}{10 g}
  \begin{cajatrampa}\footnotesize
    \begin{itemize}
    \item Contar mal los períodos (incluir el instante inicial).
    \item Modelar con una función lineal un crecimiento porcentual.
    \end{itemize}
  \end{cajatrampa}
\end{frame}

\begin{frame}[allowframebreaks]{Funciones trigonométricas seno y coseno}
  \framesubtitle{Teoría esencial}
  \begin{itemize}\small
    \item $\sin 0=0$, $\sin 90^\circ=1$; $\cos 0=1$, $\cos 90^\circ=0$.
    \item La gráfica del coseno es la del seno desplazada $90^\circ$.
    \item Un $b$ mayor comprime la gráfica: más ciclos en el mismo intervalo.
  \end{itemize}
  \begin{cajaclave}[Parámetros]\footnotesize
\begin{gather*}
y=a\sin(bx)+c \\
\text{amplitud}=|a| \\
\text{período}=\frac{2\pi}{|b|} \\
\text{recorrido}=[\,c-|a|,\ c+|a|\,]
\end{gather*}
  \end{cajaclave}
\end{frame}

\begin{frame}{Desafío}
  \framesubtitle{Funciones trigonométricas seno y coseno}
  {\small La altura de la marea, en metros, se modela con $h(t)=1{,}5\sin\!\left(\dfrac{\pi t}{6}\right)+3$, con $t$ en horas. ¿Cuál es la altura máxima?}
  \vspace{0.4em}
  \alternativa{A}{1,5 m}
  \alternativa{B}{3 m}
  \alternativa{C}{4,5 m}
  \alternativa{D}{6 m}
\end{frame}

\begin{frame}[allowframebreaks]{Solución}
  \framesubtitle{El seno vale como máximo 1.}
  {\small
  \paso{1}{Máximo: $1{,}5\cdot 1+3=4{,}5$ m.}
  \paso{2}{El período es $\dfrac{2\pi}{\pi/6}=12$ horas.}
  }
  \alternativaok{C}{4,5 m}
  \begin{cajatrampa}\footnotesize
    \begin{itemize}
    \item Confundir la amplitud con la distancia entre máximo y mínimo (que es el doble).
    \item Calcular el período como $2\pi\cdot b$ en vez de $\dfrac{2\pi}{b}$.
    \end{itemize}
  \end{cajatrampa}
\end{frame}

\begin{frame}[allowframebreaks]{Ecuación y función cuadrática (contenido de M1)}
  \framesubtitle{Teoría esencial}
  \begin{itemize}\small
    \item Ordena la ecuación en la forma $ax^2+bx+c=0$ antes de identificar coeficientes.
    \item «Una única solución» $\iff \Delta=0$.
    \item La parábola no corta el eje $X$ cuando $\Delta<0$.
  \end{itemize}
  \begin{cajaclave}[Discriminante]\footnotesize
$\Delta>0$: dos soluciones reales distintas. $\Delta=0$: una solución real (doble). $\Delta<0$: ninguna solución real.

\begin{gather*}
x=\frac{-b\pm\sqrt{\Delta}}{2a} \\
x_1+x_2=-\frac{b}{a} \\
x_1x_2=\frac{c}{a}
\end{gather*}
  \end{cajaclave}
\end{frame}

\begin{frame}{Desafío}
  \framesubtitle{Ecuación y función cuadrática (contenido de M1)}
  {\small ¿Para qué valores de $k$ la ecuación $x^2-4x+k=0$ no tiene soluciones reales?}
  \vspace{0.4em}
  \alternativa{A}{$k<4$}
  \alternativa{B}{$k=4$}
  \alternativa{C}{$k>4$}
  \alternativa{D}{$k>16$}
\end{frame}

\begin{frame}[allowframebreaks]{Solución}
  \framesubtitle{No hay soluciones reales cuando el discriminante es negativo.}
  {\small
  \paso{1}{$\Delta=(-4)^2-4\cdot 1\cdot k=16-4k$.}
  \paso{2}{$16-4k<0 \Rightarrow k>4$.}
  }
  \alternativaok{C}{$k>4$}
  \begin{cajatrampa}\footnotesize
    \begin{itemize}
    \item Olvidar el factor 4 en $4ac$.
    \item Perder una solución al dividir por $x$.
    \end{itemize}
  \end{cajatrampa}
\end{frame}

\section{Geometría}

\begin{frame}[plain]
  \vfill
  {\color{zinc500}\small Eje}\par\vspace{0.3em}{\Large\bfseries Geometría}\par\vspace{0.6em}
  {\color{zinc700} Homotecia, razones trigonométricas, circunferencia, esfera y rectas en el plano.}\par\vspace{1em}
  {\small\color{zinc500} Homotecia de figuras planas, Razones trigonométricas, Relaciones métricas en la circunferencia, Esfera: superficie y volumen, Rectas en el plano}
  \vfill
\end{frame}

\begin{frame}[allowframebreaks]{Homotecia de figuras planas}
  \framesubtitle{Teoría esencial}
  \begin{itemize}\small
    \item El signo de $k$ solo cambia el sentido, no el tamaño.
    \item Razón $-1$: equivale a una rotación de $180^\circ$.
    \item Un área nunca resulta negativa.
  \end{itemize}
  \begin{cajaclave}[Efecto de la razón]\footnotesize
Longitudes $\times|k|$, áreas $\times k^2$, volúmenes $\times|k|^3$. Con centro en el origen: $(x,y)\mapsto(kx,\,ky)$.
  \end{cajaclave}
\end{frame}

\begin{frame}{Desafío}
  \framesubtitle{Homotecia de figuras planas}
  {\small Con centro en el origen y razón $k=3$, el punto $A(2,-1)$ se transforma en:}
  \vspace{0.4em}
  \begin{tcbraster}[raster columns=2, raster column skip=0.6em, raster row skip=0.35em]
  \alternativa{A}{$(5,2)$}
  \alternativa{B}{$(6,-3)$}
  \alternativa{C}{$(-6,3)$}
  \alternativa{D}{$\left(\tfrac{2}{3},-\tfrac{1}{3}\right)$}
  \end{tcbraster}
\end{frame}

\begin{frame}[allowframebreaks]{Solución}
  \framesubtitle{Con centro en el origen, se multiplican ambas coordenadas por $k$.}
  {\small
  \paso{1}{$(3\cdot 2,\ 3\cdot(-1))=(6,-3)$.}
  }
  \alternativaok{B}{$(6,-3)$}
  \begin{cajatrampa}\footnotesize
    \begin{itemize}
    \item Multiplicar el área por $k$ en vez de $k^2$.
    \item Arrastrar el signo de $k$ a las medidas.
    \end{itemize}
  \end{cajatrampa}
\end{frame}

\begin{frame}[allowframebreaks]{Razones trigonométricas}
  \framesubtitle{Teoría esencial}
  \begin{itemize}\small
    \item Identifica qué lado es opuesto y cuál adyacente \textbf{respecto del ángulo dado}.
    \item Ningún cateto es mayor que la hipotenusa.
    \item $\tan\alpha=\frac{\operatorname{sen}\alpha}{\cos\alpha}$.
  \end{itemize}
  \begin{cajaclave}[Razones y valores notables]\footnotesize
\begin{gather*}
\operatorname{sen}\alpha=\frac{\text{op}}{\text{hip}} \\
\cos\alpha=\frac{\text{ady}}{\text{hip}} \\
\tan\alpha=\frac{\text{op}}{\text{ady}} \\
\operatorname{sen}^2\alpha+\cos^2\alpha=1
\end{gather*}

$\operatorname{sen}30^\circ=\cos 60^\circ=\frac{1}{2}$, $\operatorname{sen}60^\circ=\cos 30^\circ=\frac{\sqrt{3}}{2}$, $\tan 45^\circ=1$.
  \end{cajaclave}
\end{frame}

\begin{frame}{Desafío}
  \framesubtitle{Razones trigonométricas}
  {\small Desde un punto del suelo a 20 m de un edificio, se ve su azotea con un ángulo de elevación de $45^\circ$. ¿Cuál es la altura del edificio?}
  \vspace{0.4em}
  \alternativa{A}{10 m}
  \alternativa{B}{20 m}
  \alternativa{C}{$20\sqrt{2}$ m}
  \alternativa{D}{40 m}
\end{frame}

\begin{frame}[allowframebreaks]{Solución}
  \framesubtitle{La altura es el cateto opuesto y la distancia, el adyacente: se usa la tangente.}
  {\small
  \paso{1}{$\tan 45^\circ=\dfrac{h}{20}$.}
  \paso{2}{$h=20\cdot 1=20$ m.}
  }
  \alternativaok{B}{20 m}
  \begin{cajatrampa}\footnotesize
    \begin{itemize}
    \item Usar el coseno cuando corresponde el seno (o viceversa).
    \item Dividir en vez de multiplicar al despejar.
    \end{itemize}
  \end{cajatrampa}
\end{frame}

\begin{frame}[allowframebreaks]{Relaciones métricas en la circunferencia}
  \framesubtitle{Teoría esencial}
  \begin{itemize}\small
    \item Ángulo inscrito en una semicircunferencia: $90^\circ$.
    \item Cuadrilátero inscrito: sus ángulos opuestos suman $180^\circ$.
    \item Tangente y radio son perpendiculares en el punto de tangencia.
  \end{itemize}
  \begin{cajaclave}[Relaciones]\footnotesize
Ángulo inscrito $=\dfrac{\text{ángulo del centro}}{2}$. Longitud de arco $=\dfrac{\alpha}{360^\circ}\cdot 2\pi r$.

\begin{gather*}
\text{cuerdas: } PA\cdot PB=PC\cdot PD \\
\text{secantes: } PA\cdot PB=PC\cdot PD\ (P \text{ exterior})
\end{gather*}
  \end{cajaclave}
\end{frame}

\begin{frame}{Desafío}
  \framesubtitle{Relaciones métricas en la circunferencia}
  {\small Dos cuerdas $AB$ y $CD$ se cortan en $P$. Si $PA=4$, $PB=6$ y $PC=3$, ¿cuánto mide $PD$?}
  \vspace{0.4em}
  \alternativa{A}{2}
  \alternativa{B}{4,5}
  \alternativa{C}{8}
  \alternativa{D}{13}
\end{frame}

\begin{frame}[allowframebreaks]{Solución}
  \framesubtitle{En cuerdas que se cortan, los productos de sus segmentos son iguales.}
  {\small
  \paso{1}{$PA\cdot PB=PC\cdot PD \Rightarrow 4\cdot 6=3\cdot PD$.}
  \paso{2}{$PD=8$.}
  }
  \alternativaok{C}{8}
  \begin{cajatrampa}\footnotesize
    \begin{itemize}
    \item Tomar el ángulo inscrito igual al del centro.
    \item En secantes, multiplicar los dos segmentos exteriores en vez del exterior por la secante completa.
    \end{itemize}
  \end{cajatrampa}
\end{frame}

\begin{frame}[allowframebreaks]{Esfera: superficie y volumen}
  \framesubtitle{Teoría esencial}
  \begin{itemize}\small
    \item Duplicar el radio multiplica el volumen por 8.
    \item El área de la esfera es igual a 4 círculos máximos.
    \item Diámetro $=2\cdot$ radio.
  \end{itemize}
  \begin{cajaclave}[Fórmulas]\footnotesize
\begin{gather*}
A=4\pi r^2 \\
V=\frac{4}{3}\pi r^3 \\
V_{\text{semiesfera}}=\frac{2}{3}\pi r^3
\end{gather*}
  \end{cajaclave}
\end{frame}

\begin{frame}{Desafío}
  \framesubtitle{Esfera: superficie y volumen}
  {\small Una pelota de radio 3 cm se cambia por otra de radio 6 cm. ¿Cuántas veces mayor es el volumen de la nueva pelota?}
  \vspace{0.4em}
  \alternativa{A}{2}
  \alternativa{B}{4}
  \alternativa{C}{6}
  \alternativa{D}{8}
\end{frame}

\begin{frame}[allowframebreaks]{Solución}
  \framesubtitle{El volumen depende del cubo del radio.}
  {\small
  \paso{1}{La razón entre radios es $k=2$.}
  \paso{2}{Los volúmenes están en razón $k^3=8$.}
  }
  \alternativaok{D}{8}
  \begin{cajatrampa}\footnotesize
    \begin{itemize}
    \item Usar el diámetro como radio.
    \item Olvidar el factor $\dfrac{4}{3}$ del volumen.
    \end{itemize}
  \end{cajatrampa}
\end{frame}

\begin{frame}[allowframebreaks]{Rectas en el plano}
  \framesubtitle{Teoría esencial}
  \begin{itemize}\small
    \item Pendiente perpendicular: invierte y cambia el signo.
    \item De $ax+by+c=0$: $m=-\frac{a}{b}$.
    \item Comprueba que el punto dado satisfaga la ecuación.
  \end{itemize}
  \begin{cajaclave}[Fórmulas]\footnotesize
\begin{gather*}
d(A,B)=\sqrt{(x_2-x_1)^2+(y_2-y_1)^2} \\
M=\left(\frac{x_1+x_2}{2},\frac{y_1+y_2}{2}\right)
\end{gather*}

Paralelas: $m_1=m_2$. Perpendiculares: $m_1\cdot m_2=-1$.
  \end{cajaclave}
\end{frame}

\begin{frame}{Desafío}
  \framesubtitle{Rectas en el plano}
  {\small ¿Cuál es la distancia entre los puntos $A(1,2)$ y $B(4,6)$?}
  \vspace{0.4em}
  \alternativa{A}{3}
  \alternativa{B}{4}
  \alternativa{C}{5}
  \alternativa{D}{7}
\end{frame}

\begin{frame}[allowframebreaks]{Solución}
  \framesubtitle{Aplica la fórmula de distancia (Pitágoras con las diferencias de coordenadas).}
  {\small
  \paso{1}{$\Delta x=3$, $\Delta y=4$.}
  \paso{2}{$d=\sqrt{3^2+4^2}=\sqrt{25}=5$.}
  }
  \alternativaok{C}{5}
  \begin{cajatrampa}\footnotesize
    \begin{itemize}
    \item Solo invertir o solo cambiar el signo de la pendiente para la perpendicular.
    \item Leer mal la pendiente de la forma general.
    \end{itemize}
  \end{cajatrampa}
\end{frame}

\section{Probabilidad y estadística}

\begin{frame}[plain]
  \vfill
  {\color{zinc500}\small Eje}\par\vspace{0.3em}{\Large\bfseries Probabilidad y estadística}\par\vspace{0.6em}
  {\color{zinc700} Dispersión, probabilidad condicional, permutación y combinatoria, y modelos binomial y normal.}\par\vspace{1em}
  {\small\color{zinc500} Medidas de dispersión, Probabilidad condicional, Permutación y combinatoria, Modelos binomial y normal}
  \vfill
\end{frame}

\begin{frame}[allowframebreaks]{Medidas de dispersión}
  \framesubtitle{Teoría esencial}
  \begin{itemize}\small
    \item Igual media no implica igual dispersión.
    \item Sumar 5 puntos a todos no cambia $\sigma$.
    \item Multiplicar todos los datos por 2 duplica $\sigma$ y cuadruplica la varianza.
  \end{itemize}
  \begin{cajaclave}[Fórmulas]\footnotesize
\begin{gather*}
\sigma=\sqrt{\frac{\sum (x_i-\bar{x})^2}{n}} \\
\sigma^2=\text{varianza} \\
\text{rango}=\text{máx}-\text{mín}
\end{gather*}
  \end{cajaclave}
\end{frame}

\begin{frame}{Desafío}
  \framesubtitle{Medidas de dispersión}
  {\small Los puntajes de una prueba tienen media 60 y desviación estándar 8. Si a todos se les suman 5 puntos, ¿cuál es la nueva desviación estándar?}
  \vspace{0.4em}
  \alternativa{A}{3}
  \alternativa{B}{8}
  \alternativa{C}{13}
  \alternativa{D}{40}
\end{frame}

\begin{frame}[allowframebreaks]{Solución}
  \framesubtitle{Sumar una constante desplaza los datos sin cambiar sus distancias a la media.}
  {\small
  \paso{1}{La media pasa a 65, pero cada dato sigue a la misma distancia de ella.}
  \paso{2}{La desviación estándar se mantiene en 8.}
  }
  \alternativaok{B}{8}
  \begin{cajatrampa}\footnotesize
    \begin{itemize}
    \item Creer que sumar una constante cambia la desviación estándar.
    \item Confundir variedad de datos con dispersión.
    \end{itemize}
  \end{cajatrampa}
\end{frame}

\begin{frame}[allowframebreaks]{Probabilidad condicional}
  \framesubtitle{Teoría esencial}
  \begin{itemize}\small
    \item La condición indica el denominador.
    \item $P(A\mid B)\neq P(B\mid A)$ en general.
    \item En un árbol, se multiplica a lo largo de una rama y se suma entre ramas.
  \end{itemize}
  \begin{cajaclave}[Fórmulas]\footnotesize
\begin{gather*}
P(A\mid B)=\frac{P(A\cap B)}{P(B)} \\
P(A\cap B)=P(B)\cdot P(A\mid B)
\end{gather*}
  \end{cajaclave}
\end{frame}

\begin{frame}{Desafío}
  \framesubtitle{Probabilidad condicional}
  {\small En un curso de 40 estudiantes, 25 practican deporte; de ellos, 10 también tocan un instrumento. Si se elige al azar un estudiante que practica deporte, ¿cuál es la probabilidad de que toque un instrumento?}
  \vspace{0.4em}
  \begin{tcbraster}[raster columns=2, raster column skip=0.6em, raster row skip=0.35em]
  \alternativa{A}{$\dfrac{1}{4}$}
  \alternativa{B}{$\dfrac{2}{5}$}
  \alternativa{C}{$\dfrac{5}{8}$}
  \alternativa{D}{$\dfrac{3}{5}$}
  \end{tcbraster}
\end{frame}

\begin{frame}[allowframebreaks]{Solución}
  \framesubtitle{La condición «practica deporte» reduce el espacio muestral a 25 estudiantes.}
  {\small
  \paso{1}{Casos posibles: los 25 que practican deporte.}
  \paso{2}{Favorables: 10. $P=\dfrac{10}{25}=\dfrac{2}{5}$.}
  \paso{3}{$\frac{1}{4}=\frac{10}{40}$ usa el curso completo como espacio muestral, ignorando la condición.}
  }
  \alternativaok{B}{$\dfrac{2}{5}$}
  \begin{cajatrampa}\footnotesize
    \begin{itemize}
    \item Invertir la condición.
    \item Usar el total general como denominador cuando hay una condición.
    \end{itemize}
  \end{cajatrampa}
\end{frame}

\begin{frame}[allowframebreaks]{Permutación y combinatoria}
  \framesubtitle{Teoría esencial}
  \begin{itemize}\small
    \item ¿Importa el orden? Sí $\Rightarrow$ variación; no $\Rightarrow$ combinación.
    \item $\binom{n}{k}=\binom{n}{n-k}$.
    \item Con repetición: $n^k$.
  \end{itemize}
  \begin{cajaclave}[Fórmulas]\footnotesize
\begin{gather*}
P_n=n! \\
V^n_k=\frac{n!}{(n-k)!} \\
\binom{n}{k}=\frac{n!}{k!\,(n-k)!}
\end{gather*}
  \end{cajaclave}
\end{frame}

\begin{frame}{Desafío}
  \framesubtitle{Permutación y combinatoria}
  {\small ¿Cuántos comités de 2 personas se pueden formar con un grupo de 6 personas?}
  \vspace{0.4em}
  \alternativa{A}{12}
  \alternativa{B}{15}
  \alternativa{C}{30}
  \alternativa{D}{36}
\end{frame}

\begin{frame}[allowframebreaks]{Solución}
  \framesubtitle{En un comité no importa el orden: es una combinación.}
  {\small
  \paso{1}{$\binom{6}{2}=\dfrac{6\cdot 5}{2}=15$.}
  \paso{2}{30 contaría cada par dos veces (con orden).}
  }
  \alternativaok{B}{15}
  \begin{cajatrampa}\footnotesize
    \begin{itemize}
    \item Usar combinaciones cuando los cargos son distintos.
    \item Permitir repeticiones que el problema prohíbe.
    \end{itemize}
  \end{cajatrampa}
\end{frame}

\begin{frame}[allowframebreaks]{Modelos binomial y normal}
  \framesubtitle{Teoría esencial}
  \begin{itemize}\small
    \item $\binom{n}{k}$ cuenta en qué ensayos caen los $k$ éxitos.
    \item La normal es simétrica: la mitad de los datos está bajo la media.
    \item Un $z$ negativo indica un dato bajo la media.
  \end{itemize}
  \begin{cajaclave}[Fórmulas]\footnotesize
\begin{gather*}
P(X=k)=\binom{n}{k}p^k(1-p)^{n-k} \\
E(X)=np \\
z=\frac{x-\mu}{\sigma}
\end{gather*}

$P(\mu-\sigma\le X\le\mu+\sigma)\approx 0{,}68$; $P(\mu-2\sigma\le X\le\mu+2\sigma)\approx 0{,}95$.
  \end{cajaclave}
\end{frame}

\begin{frame}{Desafío}
  \framesubtitle{Modelos binomial y normal}
  {\small Los puntajes de un test siguen una distribución normal de media 500 y desviación estándar 100. ¿Qué porcentaje aproximado de los puntajes está entre 400 y 600?}
  \vspace{0.4em}
  \alternativa{A}{34 \%}
  \alternativa{B}{50 \%}
  \alternativa{C}{68 \%}
  \alternativa{D}{95 \%}
\end{frame}

\begin{frame}[allowframebreaks]{Solución}
  \framesubtitle{400 y 600 están a una desviación estándar de la media.}
  {\small
  \paso{1}{$z=\dfrac{400-500}{100}=-1$ y $z=\dfrac{600-500}{100}=1$.}
  \paso{2}{Entre $z=-1$ y $z=1$ queda cerca del 68 \%.}
  }
  \alternativaok{C}{68 \%}
  \begin{cajatrampa}\footnotesize
    \begin{itemize}
    \item Olvidar el coeficiente binomial.
    \item Leer el 95 \% como «a una desviación de la media» en vez de a dos.
    \end{itemize}
  \end{cajatrampa}
\end{frame}

\section{Habilidades transversales}

\begin{frame}[plain]
  \vfill
  {\color{zinc500}\small Eje}\par\vspace{0.3em}{\Large\bfseries Habilidades transversales}\par\vspace{0.6em}
  {\color{zinc700} Formatos de pregunta que exigen razonamiento más que cálculo.}\par\vspace{1em}
  {\small\color{zinc500} Evaluación de suficiencia de datos}
  \vfill
\end{frame}

\begin{frame}[allowframebreaks]{Evaluación de suficiencia de datos}
  \framesubtitle{Teoría esencial}
  \begin{itemize}\small
    \item Suficiente = determina un valor único.
    \item No mezcles la información de (1) al analizar (2).
    \item «Ambas juntas» solo cuando ninguna basta sola.
  \end{itemize}
  \begin{cajaclave}[Orden de análisis]\footnotesize
1) Analiza (1) sola. 2) Analiza (2) sola. 3) Solo si ninguna basta, analiza ambas juntas.

Si las dos bastan por separado, la respuesta es «cada una por sí sola».
  \end{cajaclave}
\end{frame}

\begin{frame}{Desafío}
  \framesubtitle{Evaluación de suficiencia de datos}
  {\small Se puede determinar el valor de $x$ si:

(1) $x^2=9$.

(2) $x>0$.}
  \vspace{0.4em}
  \alternativa{A}{(1) por sí sola}
  \alternativa{B}{(2) por sí sola}
  \alternativa{C}{Ambas juntas, (1) y (2)}
  \alternativa{D}{Se requiere información adicional}
\end{frame}

\begin{frame}[allowframebreaks]{Solución}
  \framesubtitle{Con (1) hay dos valores posibles; con (2) infinitos; juntas, uno solo.}
  {\small
  \paso{1}{(1): $x=3$ o $x=-3$. No es único.}
  \paso{2}{(2): cualquier positivo. No basta.}
  \paso{3}{Juntas: $x=3$. Son suficientes ambas juntas.}
  }
  \alternativaok{C}{Ambas juntas, (1) y (2)}
  \begin{cajatrampa}\footnotesize
    \begin{itemize}
    \item Usar sin darse cuenta datos de (1) al evaluar (2).
    \item Elegir «ambas juntas» cuando una sola ya era suficiente.
    \end{itemize}
  \end{cajatrampa}
\end{frame}
%%</CONTENIDO>%%

\end{document}
